\zlabel{triptych:concise:commentary:start}
\section{The Propers: Detailed Commentary}

\subsection{How can rescued people still need rescue?}

The complete Psalm 43(44) makes the Gradual's praise difficult. It begins by
recalling the ancestors' possession of the land through God's help. It then
praises rescue and promises further praise, before lamenting defeat and
dispersion. The speakers insist that they have remained faithful. Their
suffering therefore cannot simply be explained as their having abandoned
the God whom the Psalm thanks.

Augustine's answer begins with the change of tense. The psalmist speaks
of rescue as something done, then promises praise still to come.
Augustine reads the rescue as future and compares the words about pierced
hands and feet in Psalm 21(22). That psalm speaks in the past tense
of Christ's Passion, still future for the prophet. God's foreknowledge
allows the prophet the same certainty about the future as about an event
already witnessed. The lament describes the Church's suffering between the
ancestors' deliverance and the everlasting kingdom.\footnote{Augustine,
\work{Enarrationes in Psalmos}, Ps 43, Latin §§8--9, NPNF Ps XLIV §§6--7
for the verbs and Passion comparison, and §§4, 8, 13 for the lament and
the interval.}

Bellarmine takes the past tense in its ordinary sense and places the lament
at a later affliction.\footnote{Bellarmine,
Psalm 43, native vv. 7--8 (Vulgate vv. 8--9), followed by the onset of
the lament at native v. 9 (Vulgate v. 10).}

The Alleluia asks to be heard just after the Gradual sings of rescue.
In Augustine's account, the Church can praise a rescue it still awaits.
In Bellarmine's, it can thank God for earlier rescues while asking him
for help again.

The Introit's Psalm 84(85) raises a similar difficulty within its own
verses. After speaking of captivity ended, the Psalm asks God to turn
his people again. Augustine asks why the prophet would pray for an act
he had just described as accomplished. He concludes that the prophet
foresees the restoration and asks for what he sees to come about.
Augustine explains Jacob's captivity through Romans 7, where Paul
describes being captive to the law of sin. He also takes up mortality
when the psalmist asks for God's anger to end. For Christians giving
thanks with Paul for forgiveness already received, freedom from
mortality still lies ahead.\footnote{Augustine, \work{Enarrationes in
Psalmos}, Ps 84, Latin §6 on the prophet's prayer, §§4--5 on sin and
forgiveness, and §7 on mortality. Col 1:12--14.}

The Epistle supplies a related distinction within one passage. Chrysostom
explains rescue from darkness and transfer into the Son's kingdom as a
double benefit. Escape from an oppressor would already be a great gift;
admission into the beloved Son's kingdom gives a new belonging. Paul's
continuing prayer concerns people who have truly received this benefit.
It asks that their knowledge and life grow rather than treating them as
people to whom nothing has happened.\footnote{Chrysostom,
\work{Homilies on Colossians} 2, on Col 1:9--14.}

The second interpretation asks a different question about the same rescue:
what becomes of the person God brings into the kingdom? Aquinas attends to
Paul's assertion that the Father made the recipients worthy. Worthiness is
itself God's effect in them. It cannot be an independent prior achievement
explaining why they received the first grace. Chrysostom makes the point
with the example of a king granting office. An honour given to an incapable
person would leave him exposed to ridicule. God's gift includes the fitness
to receive it.\footnote{Aquinas, \work{Super Colossenses} 1, lectio 3,
on Col 1:12, native anchor 87848; Chrysostom, \work{Homilies on Colossians}
2.}

Paul asks \textit{ut ambuletis digne Deo}, that his hearers may walk
worthily of God. They are to bear fruit in good
works. Aquinas explains why Paul asks for growth in knowledge as well
as good works. Acting on what the faithful have learned disposes them
to understand God's will more deeply. In the first interpretation,
such work continues while the faithful await final deliverance. In
the second, God makes the faithful capable of the work for which Paul
prays, just as the Collect asks him to stir their wills so they may
bear fruit.\footnote{1962 Missal, p. 423, nos. 1770--1771; Aquinas,
\work{Super Colossenses} 1, lectio 3, on Col 1:9--12, native anchor
87848.}

\subsection{Why do peace and faith require work, flight and discernment?}

Jeremiah's letter gives the first difficulty an everyday form. The
deportees must plant and form households in the city from which they hope
to return. Their hope does not cancel those duties. Jerome explains the
promise against the false prophets' expectation of immediate restoration.
He also relates the future rest spiritually to the troubles and dispersion
of this life. In his next unit Jerome explains the gathering clause
that the Introit keeps as a promise about the historical exiles.\footnote{Jer
29:4--14; Jerome, \work{In Ieremiam} V,
\textit{Vers. 10 seqq.} on Jer 29:10--14a and opening of
\textit{Vers. 14 seqq.} on Jer 29:14b, PL 24, cols. 893--894
(facsimile PDF p. 451).}

Jesus' warning to flee concerns a different situation and requires a
different action. In Judea a person must leave the roof or field without
returning for possessions. Chrysostom explains pregnancy and nursing through
the difficulty of escape. These circumstances reveal vulnerability; they
are no evidence that the afflicted have failed to trust. The command to
flee requires those who believe Jesus to use the warning he gives.

Chrysostom distinguishes that catastrophe from the public final coming.
He explains ``then'' in verse 23 by comparing Matthew's move from the death
of Herod to John the Baptist's preaching. Matthew introduces the preaching
with ``in those days'' and leaves the intervening years unmentioned.
Chrysostom understands ``then'' as leaving a similar interval after the
siege.\footnote{Chrysostom,
\work{Homilies on Matthew} 76.2.}

Chrysostom hears ``immediately'' in verse 29 as a promise that Christ comes
as soon as the final trial ends. Jesus has just warned about deception,
and Chrysostom understands that trial as the work of Antichrist and the
false prophets. He expects God to limit the final trial for the elect's
sake, as he limited the siege.\footnote{Chrysostom, \work{Homilies on Matthew} 76.2--3,
on Mt 24:23--29.}

An account confined to Jerusalem readily explains the roof, field and
season. It does not exhaust the universal manifestation and gathering as
Chrysostom understands them. An account concerned only with a future
catastrophe loses the immediate force of the Judean instructions.
The enduring ``generation'' must be understood accordingly. Chrysostom
identifies it by faith and practice: the generation of believers remains
through the trials. Aquinas also distinguishes this explanation from a
contemporary cohort appropriate to a restricted Jerusalem reference.
The broader promise concerns a people whose continuance is not measured
by the lifespan of its first members.\footnote{Chrysostom,
\work{Homilies on Matthew} 77; Aquinas, \work{Super Matthaeum} 24,
lectio 3, on Mt 24:34.}

The first interpretation finds in Paul's prayer two kinds of endurance.
Aquinas distinguishes patience under adversity from longanimity under
delay. A person may bear the immediate difficulty yet tire because the
reward remains absent. Chrysostom finds reason for joy in the gifts God
has already given. He also recalls that Job blessed God for what he had
given while grieving its loss.\footnote{Aquinas, \work{Super Colossenses}
1, lectio 3, on Col 1:11; Chrysostom, \work{Homilies on Colossians} 2,
on the transition from 1:11 to 1:12.}

The second asks how desire can make a person credulous. A spectacular sign
may attract someone who wants relief, although its claimant contradicts
Christ. Augustine applies the contrast between private rooms and lightning
to the Church's public presence. He hears the claim that Christ is here
or there as the claim of a schism to possess him in one place, while
east and west describe the reach of the Gospel from Jerusalem to the
whole world. Chrysostom applies the lightning to Christ's final
appearance to all at once.\footnote{Augustine, \work{Quaestiones
euangeliorum} I.38, on Mt 24:23--27, registered Latin transcription,
line 43; Chrysostom, \work{Homilies on Matthew} 76.2--3.}

Aquinas's discussion of deception reports Origen's distinction between
human vulnerability and the preserving effect of divine election.
Aquinas also gives the alternative of people elect only in appearance.
In either case, Jesus warns people who believe they belong to him to
heed what he has already told them.\footnote{Aquinas,
\work{Super Matthaeum} 24, lectio 3,
on Mt 24:24, native anchor 87397. Origen's explanation is reported
through Aquinas.}

The body and eagles in verse 28 picture a gathering. Chrysostom
identifies the gathered multitude as angels, martyrs and saints.
Aquinas relates the body to Christ's Passion and the eagles to the
saints drawn to him. In the
first interpretation, the images give the dispersed people a destination
when Christ sends his angels to gather the elect. In the second,
Aquinas's account directs desire toward the Christ who gave himself in
the Passion, while Augustine's account of the lightning gives his
people a way to assess claims to exclusive access.\footnote{Chrysostom,
\work{Homilies on Matthew} 76.3;
Aquinas, \work{Super Matthaeum} 24, lectio 3, on Mt 24:28, 30--31.}

Related to the Secret, the warning makes attachment concrete. A coat can
be a legitimate possession and still be something one must leave when
Christ commands flight. The exiles must work; those fleeing must depart.
Conversion of desire means responding to God's actual command rather
than making either possession or immediate escape an unconditional good.

\subsection{Why ask for a gift God already intends to give?}

Jerome raises this question directly over Jeremiah's promise. If return
is assured by God, why does the passage still require invocation and
prayer? His answer concerns the recipients: prayer prepares them to receive
the promised benefit. The Lord's generosity calls forth a real response.
The Introit abbreviates the letter, but its context still requires
whole-hearted seeking.

Aquinas distinguishes two acts in Paul's words. Prayer raises the mind
to God and prepares the person to ask. In petition the person asks for
particular gifts. Here Paul asks for knowledge of God's will, the practice
of virtue and endurance under adversity.\footnote{Jerome, \work{In Ieremiam} V,
\textit{Vers. 10 seqq.} on Jer 29:10--14a, PL 24, cols. 893--894;
Aquinas, \work{Super Colossenses} 1, lectio 3, native anchor 87848,
on praying and petitioning.}

We can apply Aquinas's distinction to the Secret. The Church asks God to
receive the people's offerings and prayers and turn their hearts toward
himself. Fromage
explains this prayer as a passage from earthly attachment to heavenly
desire. The historical English's ``enjoyments of heaven'' renders
a Latin phrase that names what the offerers are to come to
desire.\footnote{Fromage, vol. XI, p. 513, comment on the Secret;
1962 Missal, p. 424, no. 1776; Lasance/Walsh, Last Sunday after Pentecost,
Secret, p. 740.}

The Alleluia asks God to hear the cry from the depths before the Gospel.
The Offertory antiphon repeats its opening words. We can hear the cry as
a plea for the change of heart the Church asks for in the Secret.
Augustine explains why the cry already shows
hope: someone who calls out still seeks release. Bellarmine compares the
sinner to someone crying from far below. He locates the distance in the
sinner's dissimilarity to God, who remains present by essence and power.
The later verses of the Psalm explain why penitence needs both knowledge
of misery and confidence in mercy. Bellarmine says that whoever fails
to recognise his misery will decline medicine, while whoever fails to
recognise God's mercy despairs of cure.\footnote{Augustine,
\work{Enarrationes in Psalmos}, Ps 129 §§1--2; Bellarmine, Psalm 129,
on vv. 1--2, 4--5; 1962 Missal, pp. 423--424, nos. 1773, 1775.}

For the interpretation of deliverance begun, this explains why the
redeemed keep praying before their suffering ends. For the interpretation
of the changed petitioner, it explains why confidence includes
acknowledgment of one's faults. Those are complementary answers, but
they cannot replace each other. Suffering does not by itself prove a
fault in its sufferer, while repentance concerns a defect that really
needs cure.

In the Preface the worshippers thank the Father. Paul asks the faithful
to thank him for bringing them into his Son's kingdom. Chrysostom explains
that forgiveness in the Son makes the transfer possible. Aquinas, commenting
on Col 1:15--17 beyond the appointed lesson, explains the Son's possession
of the same invisible divine nature. All things are created through him and
sustained in him.\footnote{Chrysostom,
\work{Homilies on Colossians} 2, on 1:13--14; Aquinas,
\work{Super Colossenses} 1, lectio 4.}

The first interpretation understands the Sanctus as sharing now in
the praise the Church hopes to offer without end. The Preface speaks
of the angels praising God ``with one voice,'' and the worshippers
join their ``Holy, holy, holy, Lord God of hosts.'' The second
interpretation relates the Secret's heavenly desire to the God whom the worshippers
praise.\footnote{1962 Missal, p. 293, no. 1082; Lasance/Walsh,
Most Holy Trinity Preface and Sanctus, pp. 773--774.}

\subsection{What does the Communion's assurance allow us to expect?}

Mark places the promise after the withered fig tree and Jesus' command
to have faith in God. The Communion keeps the instruction to believe
we will receive what we ask. Its ending in the historical Lasance/Walsh
form is ``it shall be done unto you.'' Challoner's Bible reads
``they shall come unto you.'' In the next verse Jesus tells those who
stand in prayer to forgive so that the Father will forgive them.\footnote{Mk 11:12--26, especially 22--25;
1962 Missal, p. 424, no. 1777; Lasance/Walsh, Last Sunday after
Pentecost, Communion, p. 740.}

Theophylact explains the appointed verse itself. Christ promises
receipt to the person who asks with faith. Theophylact describes this
as giving oneself wholly to God, with the image of someone praying
in tears as though clasping the Lord's feet. In the following verse
he understands forgiving a person who has caused offence as another
way to receive what one asks.\footnote{Theophylact of Ohrid,
\work{In Marcum}, Basel 1525, fols. 69r--v, on Mk 11:24--25.}

Bede explains the surrounding sayings about faith, the mountain
and forgiveness. In the mountain he sees something God can move when
necessity requires. He recounts how Gregory of Neocaesarea prayed at
night over a church site hemmed in by the sea and a mountain. The
next morning Gregory found that the mountain had yielded enough ground
for the builders.\footnote{Bede, \work{In Marci evangelium expositio} III,
on Mk 11:22--23, registered Latin transcription, line 338.}

He also interprets the mountain spiritually as the devil's pride.
In explaining the command to forgive, he describes effective faith
as faith working through love. He then addresses those who have
not reached the perfection of commanding spiritual mountains.
They should ask pardon, and he assures them of receiving it if
they first forgive those who have sinned against them.\footnote{Bede,
\work{In Marci evangelium expositio} III, on Mk 11:25, registered
Latin transcription, line 339.}

Bede gives a specific assurance about pardon. His account of Gregory
shows the kind of help God can give when a need calls for it. Aquinas
explains how prayer prepares the person who asks. Theophylact's account
of wholehearted trust helps us read the Communion alongside the Secret.
A person seeking relief in a conflict can ask God for the capacity to
forgive. The gift of that capacity answers the prayer for a changed
heart even while the conflict continues.

The first interpretation looks to the Gospel's command to pray about
the time of the flight. Jesus tells the disciples to ask for help in a
real danger and promises to shorten the days for the elect. The very
people for whose sake he shortens the trial still have to face it.
Continued danger cannot tell us whether another person has prayed with
faith.\footnote{Mt 24:15--22.}

The two moral consequences remain distinct. The person facing
adversity needs endurance and truthful hope. The person asking
pardon while refusing to forgive needs conversion. An assurance
of future fulfilment alone leaves the latter demand unanswered;
an instruction to correct desires alone fails to explain the
innocent sufferer's waiting.

\subsection{What remains to be healed after Communion?}

The Postcommunion speaks after reception and asks the mysteries'
medicinal gift to cure what remains faulty in the mind. Its confidence
is expressed through a further request. Fromage explains the actual
prayer as the cure of remaining defects. Augustine describes
sacramental food and drink as received in memory of the Passion
for spiritual health through God's invisible operation. Augustine's
account helps explain how the mysteries already received can still
be the means of the healing for which the communicants ask.\footnote{Fromage,
vol. XI, p. 513, comment on the Postcommunion; Augustine,
\work{De Trinitate} III.4.10.}

In the interpretation of deliverance begun, the remaining need extends
to the final overcoming of mortality. Augustine's exposition of the
whole Psalm 129 understands its morning watch through Christ's
Resurrection. Those later verses are outside the two chant
appointments, but they explain his hope: the Head has risen in the
flesh his members share, and they await their own resurrection.
Forgiveness received now does not mean that they have already
risen as he has.\footnote{Augustine, \work{Enarrationes} 129.3--8.}

This is that interpretation's allegorical sense: captivity and
deliverance are fulfilled through Christ and his body. Its
anagogical sense concerns the final gathering and resurrection.
These fulfilments follow from the concrete needs already heard:
deportees await return, the Psalm's faithful suffer, and the
disciples must flee. The spiritual senses do not remove those
literal situations.

In the interpretation of the changed petitioner, the immediate
need concerns the recipient's continuing faults. The moral request
is to accept cure, including the conversion of desire and the
forgiveness of another person. Its allegorical sense concerns the
Church receiving and offering in Christ, whose Blood redeems the
people brought into his kingdom. Its anagogical hope is the
saints' inheritance and full communion with God. Aquinas identifies
God himself as the saints' portion when explaining Col 1:12.
Heavenly desire thus seeks the giver with whom the healed person
is to live.\footnote{Aquinas, \work{Super Colossenses} 1, lectio 3,
on the saints' portion.}
